\documentclass[12pt,openany]{book}
%控制页边距，平衡打印和阅读
\usepackage{geometry}
\geometry{
	a4paper,
	top=25.4mm, bottom=25.4mm,
	left=20mm, right=20mm,
	headheight=2.17cm,
	headsep=4mm,
	footskip=12mm
}
% 调整行距与显示
\usepackage{microtype}
\usepackage{setspace}
\setstretch{1.3} 
\usepackage[fontset=none]{ctex}
\ctexset{
	fontset = none,
	punct = quanjiao, 
}
% 中文字体（全部开源）
\setCJKmainfont{LXGWWenKaiGB}[
UprightFeatures = {Font = * Regular},
BoldFeatures    = {Font = * Medium},
]
\setCJKsansfont{Noto Sans CJK SC}[
UprightFeatures = {Font = * Regular},
BoldFeatures    = {Font = * Bold},
]

\setCJKmonofont{LXGWWenKaiMonoGB}[
UprightFeatures = {Font = * Regular},
BoldFeatures    = {Font = * Medium},
]
% 英文字体（全部开源，Noto系列与LXGW风格统一）
\usepackage[math-style=TeX]{unicode-math}
\setmainfont{STIX Two Text}[Ligatures=TeX]
\setsansfont{Noto Sans}[Ligatures=TeX]
\setmonofont{Fira Mono}[Ligatures=TeX]
\setmathfont{STIX Two Math}
%定义非全局默认的字体
\newCJKfontfamily\mengshen{Mengshen-Handwritten}
\newCJKfontfamily\youzai{Yozai}[
UprightFeatures = {Font = * Regular},
BoldFeatures    = {Font = * Medium}
]
\newCJKfontfamily\sisong{Source Han Serif SC}[
UprightFeatures = {Font = * Regular},
BoldFeatures    = {Font = * Heavy},
]
%补充字符命令
\usepackage[normalem]{ulem}
\usepackage{xcolor}%避免重复加载xcolor
%数学字符命令: 花体与符号
\usepackage{mathrsfs,mathtools,esint}
\let\eth\relax
\let\digamma\relax
\let\backepsilon\relax
\let\smallsetminus\relax
\usepackage{amssymb}
%调整页眉页脚
\usepackage{titlesec, titletoc}				
\usepackage{fancyhdr}
\fancyhf{}
\definecolor{structurecolor}{RGB}{122,122,142}
\renewcommand{\headrule}{\color{structurecolor}\hrule width\textwidth}
\pagestyle{fancy}
\renewcommand{\headrulewidth}{1pt}
\fancypagestyle{plain}{\renewcommand{\headrulewidth}{0pt}\fancyhf{}\renewcommand{\headrule}{}}
\fancyhead[c]{\color{structurecolor}\youzai\rightmark}
\fancyfoot[c]{\color{structurecolor}\small\thepage}
%设置章\节\小节等形式
\titleformat{\chapter}[display]{\sisong}
{\color{black}\filleft
	\parbox{1cm}{\vbox to 1.5cm{\vfill\hbox to 4cm{\hfill\Huge \bfseries \color{black}{Chapter} \thechapter \hfill}}}}
{1ex}
{\color{black} \titlerule[2pt]\LARGE\bfseries \filright \vspace*{0.5em}}
[\vspace*{0.5em} {\titlerule[2pt]}]
\titleformat{\section}[frame]{\normalfont\color{black}}{\footnotesize \enspace \large \textcolor{black}{\S \,\thesection}\enspace}{6pt}{\Large\filcenter \bfseries \sisong }
\titleformat{\subsection}[hang]{\bfseries\youzai }{\large\bfseries\color{black}\thesubsection\enspace}{1pt}{\color{black}\LARGE\bfseries\filright}
\titleformat{\subsubsection}[hang]{\bfseries\youzai}{\large\bfseries\color{black}\thesubsubsection\enspace}{1pt}{\color{black}\large\bfseries\filright}
%数学专用环境: 公理\定义\定理\引理\命题\性质\例题\解\证明
\usepackage{amsthm}
\usepackage[most]{tcolorbox}
% 定义标题样式：标题用思源宋体，正文用文楷 
\newtheoremstyle{plainstyle}
{3pt}{3pt}{\normalfont}{}{\bfseries\sisong}{.}{.5em}{}
\theoremstyle{plainstyle}
%普通环境,不做特殊强调
\newtheorem{theorem}{定理}[chapter]
\newtheorem{definition}[theorem]{定义}
\newtheorem{lemma}[theorem]{引理}
\newtheorem{proposition}[theorem]{命题}
\newtheorem{property}[theorem]{性质}
\newtheorem{example}[theorem]{例题}
% 特殊环境:公理,定理,命题
\definecolor{axicolor}{RGB}{255, 250, 210}
\definecolor{axiframe}{RGB}{200, 150, 50}
\newtcbtheorem[use counter=theorem]{axiom}{公理}
{
	enhanced, 
	breakable,                    
	colback=axicolor,           
	colframe=axiframe,     
	coltitle=white,              
	fonttitle=\bfseries,         
	attach boxed title to top left={xshift=5mm, yshift=-2mm},
	boxed title style={
		size=small,
		colback=axiframe,
		colframe=axiframe,
	},
	separator sign none,         
	before upper={\parindent2em} 
}{axi}
\definecolor{thmcolor}{RGB}{220, 250, 230}
\definecolor{thmframe}{RGB}{100, 180, 130}
\newtcbtheorem[use counter=theorem]{keytheorem}{定理}
{
	enhanced, 
	breakable,                    
	colback=thmcolor,           
	colframe=thmframe,     
	coltitle=white,              
	fonttitle=\bfseries,         
	attach boxed title to top left={xshift=5mm, yshift=-2mm},
	boxed title style={
		size=small,
		colback=thmframe,
		colframe=thmframe,
	},
	separator sign none,         
	before upper={\parindent2em} 
}{thm}
\definecolor{procolor}{RGB}{225, 245, 250}
\definecolor{proframe}{RGB}{100, 160, 180}
\newtcbtheorem[use counter=theorem]{keyproposition}{命题}
{
	enhanced, 
	breakable,                    
	colback=procolor,           
	colframe=proframe,     
	coltitle=white,              
	fonttitle=\bfseries,         
	attach boxed title to top left={xshift=5mm, yshift=-2mm},
	boxed title style={
		size=small,
		colback=proframe,
		colframe=proframe,
	},
	separator sign none,         
	before upper={\parindent2em} 
}{pro}
%证明和解
\renewenvironment{proof}[1][证明]{\par\underline{\textbf{\sisong #1.}} \;}{\qed\par}
\newenvironment{solution}[1][解]{\par\underline{\textbf{\sisong #1.}} \;}{\qed\par}
%引用与交叉引用
\usepackage[unicode]{hyperref} 
%注释与分栏
\usepackage{footnote}
\usepackage{multicol}
%表格处理方案
\usepackage{tabularray}
\UseTblrLibrary{booktabs}
%完整的图片处理方案
\usepackage{graphicx}
\graphicspath{{./picture/}} 
\usepackage{caption}
\captionsetup[figure]{font=small, skip=5pt}
\usepackage{float}
\usepackage{subfig}
\usepackage{wrapfig}
\usepackage{chngcntr}
\counterwithin{figure}{section}
\usepackage{tikz}
\usetikzlibrary{angles, arrows.meta, quotes, calc}
\usepackage{pgfplots}
\pgfplotsset{compat=1.16}
\usetikzlibrary{arrows.meta, positioning, calc, intersections}
\usepackage{tikz-cd}
%留言\插入\文艺
\definecolor{notecolor}{RGB}{255, 240, 245}
\newtcolorbox{note}[1][]{
	enhanced,
	breakable,
	colback=notecolor,
	colframe=red!70!black,
	coltitle=white,
	fonttitle=\bfseries,
	title=注意,
	boxed title style={
		size=small,
		colback=red!70!black,
		colframe=red!70!black
	},
	overlay unbroken and first={
		\ifx\relax#1\relax\else
		\node[anchor=north east, font=\scriptsize] 
		at (frame.north west) {#1};
		\fi
	},
	left=5mm,
	right=5mm,
	top=3mm,
	bottom=3mm,
	before upper={\parindent2em}
}
\definecolor{joycolor}{RGB}{235, 225, 250}
\definecolor{joyframe}{RGB}{150, 120, 180}
\newtcolorbox{joy}{
	fontupper=\mengshen,
	colback=joycolor,
	colframe=joyframe,
	arc=3pt,
	boxrule=0.5pt,
	left=6pt,
	right=6pt,
	top=4pt,
	bottom=4pt
}
\newtcolorbox{sad}{
	fontupper=\sisong,
	enhanced,
	breakable,
	sharp corners,
	boxrule=0.8pt,
	colback=structurecolor!5,
	colframe=structurecolor,
	left=6pt,
	right=6pt,
	top=4pt,
	bottom=4pt
}
%支持代码
\usepackage{minted}
\setminted{
	breaklines=true,
	frame=single,
	framesep=5pt,
	fontsize=\small,
	linenos=true,
	numbersep=5pt,
	bgcolor=notecolor!30,
	style=vs
}
% 设置封面
\usepackage{titling}
\renewcommand*{\maketitle}{
	\begin{titlepage}
		\newgeometry{margin = 0in}
		\parindent=0pt
		\includegraphics[width=\linewidth]{芙宁娜}
		\vfill
		\begin{center}
			\parbox{0.618\textwidth}{
				\hfill {\bfseries\sisong\Huge \thetitle} \\[0.6pt]  
				\rule{0.618\textwidth}{4pt} \\ 
			}
		\end{center}
		\vfill
		\begin{center}
			\parbox{0.618\textwidth}{
				\hfill\Large
				\begin{tabular}{r|}
					作者: \theauthor \\ 
					时间: \thedate \\
				\end{tabular}
			}
		\end{center}
		\vfill
		\begin{center}
			\parbox[t]{0.7\textwidth}{\centering  }
		\end{center}
		\vfill
	\end{titlepage}
	\restoregeometry
	\thispagestyle{empty}
}